\documentclass{article}

\usepackage{PRIMEarxiv}
\usepackage[utf8]{inputenc}
\usepackage[T1]{fontenc}
\usepackage{hyperref}
\usepackage{url}
\usepackage{booktabs}
\usepackage{amsfonts}
\usepackage{microtype}
\usepackage{graphicx}
\graphicspath{{media/}}

\title{When KV Cache Meets Backend Regimes: \\ Non-Monotonic Latency and Speedup Collapse on Apple MPS}

\author{
  Willy Fitra Hendria \\
  Independent Researcher \\
  Seoul, South Korea \\
  \texttt{willyfitrahendria@gmail.com}
}

\begin{document}
\maketitle

\begin{abstract}
Autoregressive inference is generally assumed to scale smoothly with sequence length, and key--value (KV) caching is widely treated as a universally beneficial optimization for accelerating decoding. In this work, we challenge both assumptions through a systems study of transformer inference on Apple's Metal Performance Shaders (MPS) backend. Using GPT-2-family models deployed via LitServe, a production-grade inference serving framework, we observe a systematic deviation from expected behavior: latency does not increase monotonically with sequence length, but instead exhibits abrupt latency spikes (up to $21\times$ above the pre-threshold baseline) within specific token ranges before recovering at longer lengths.

Through controlled experiments and token-level profiling, we find evidence consistent with the interpretation that this non-monotonic behavior is not explained by memory pressure or prefill cost, and is tightly linked to backend execution dynamics under cache-enabled decoding. Importantly, KV cache remains \emph{nominally} faster than cache-off across all tested lengths, but its practical advantage collapses sharply at specific sequence lengths (e.g., from $4.9$--$20.3\times$ down to $1.9\times$ at 512 tokens in our rerun). In this regime, the practical benefit of KV caching is largely neutralized despite its lower theoretical complexity. Disabling KV cache changes the latency profile and weakens the abrupt latency spike at 512 tokens, but does not restore monotonic scaling in our measurements. These results indicate a backend-sensitive performance fragility: KV caching is not the sole cause of the anomaly, but appears to interact with backend execution regime boundaries: when KV tensor growth crosses an internal threshold, the interaction produces a localized and severe performance collapse. Further analysis reveals an asymmetric mechanism under cache-enabled decoding: entry into the anomalous regime correlates with total sequence length at runtime, while recovery correlates with the maximum generation budget specified at initialization.

These findings are consistent with the hypothesis of discrete execution regimes in the MPS backend, where kernel selection or memory access patterns change abruptly as KV cache tensors grow. Our results highlight the importance of hardware- and backend-aware evaluation for autoregressive inference, and caution against relying on coarse or short-context benchmarks when assessing real-world performance.
\end{abstract}

\keywords{Apple Silicon \and MPS \and KV cache \and non-monotonic latency \and execution regimes \and autoregressive inference}

\section{Introduction}

Autoregressive inference is commonly assumed to scale smoothly with sequence length, and key--value (KV) caching~\cite{kwon2023pagedattention} is widely treated as a universally beneficial optimization. On Apple Silicon, where the PyTorch MPS backend~\cite{pytorch2022mps} enables GPU-accelerated inference as an alternative to CPU execution, prior evaluations focus on identifying when accelerators outperform CPUs~\cite{alizadeh2024llmflash}, typically reporting smooth, monotonically increasing latency trends across model sizes and sequence lengths.

These assumptions break down for autoregressive decoding. Unlike standard batch inference, decoding proceeds token by token, and the KV cache accumulates state with each generated token, creating evolving tensor shapes and memory access patterns that coarse throughput measurements do not capture.

This paper begins with a standard question: \textit{when does MPS outperform CPU?} While we reproduce a model-size-dependent crossover, deeper analysis reveals that this question is insufficient.

Instead, we uncover a more critical phenomenon:

\begin{quote}
\textbf{Autoregressive decoding on MPS exhibits non-monotonic latency scaling, with severe instability at specific sequence lengths when KV cache is enabled.}
\end{quote}

We demonstrate through controlled experiments that (1) the instability is not due to model size or memory limits, (2) it emerges during decoding rather than prefill, and (3) KV cache remains nominally faster at all tested lengths, but its practical advantage is largely neutralized at pathological lengths (speedup collapses from $4.9$--$20.3\times$ to $1.9\times$), while cache-off runs still show residual non-monotonicity.

These findings suggest that backend execution dynamics can dominate algorithmic complexity in determining autoregressive inference performance at scale. We further observe that entry into and recovery from the anomalous regime follow different signals (total sequence length for onset versus the static generation budget for recovery), consistent with an asymmetric onset/recovery dependence in backend execution behavior.

\paragraph{Contributions.}
\begin{itemize}
\item Discovery of non-monotonic latency instability in MPS autoregressive decoding, with the strongest spikes under KV cache enabled (up to $21\times$ above the pre-threshold baseline),
\item Controlled ablation evidence of a backend-sensitive KV-cache trade-off: substantial speedup at most lengths, but severe speedup collapse at pathological lengths, while cache-off runs retain residual non-monotonic behavior,
\item Characterization of an asymmetric onset/recovery pattern: onset correlates with total sequence length at runtime, while recovery correlates with the static \texttt{max\_tokens} budget at initialization,
\item A zero-model-change mitigation: tuning \texttt{max\_tokens} away from the observed anomalous region can avoid severe latency spikes in our setup.
\end{itemize}

\section{Related Work}

\subsection{LLM Inference Serving Systems}

Efficient serving of transformer-based language models has been extensively studied, primarily on CUDA-based GPU hardware.
ORCA~\cite{yu2022orca} introduces iteration-level scheduling and continuous batching to maximize GPU utilization across variable-length requests.
vLLM~\cite{kwon2023pagedattention} proposes PagedAttention, which manages KV cache memory using a virtual-memory paging abstraction to reduce fragmentation and enable larger batch sizes, achieving 2--4$\times$ throughput improvement over ORCA on standard workloads.
Both systems treat KV cache as a universally beneficial optimization and focus on maximizing throughput under high request load.
Our work differs in studying a fundamentally different failure mode: not memory fragmentation or scheduling inefficiency, but pronounced non-monotonic latency behavior suggestive of backend-dependent execution path changes under KV cache growth on a non-CUDA backend.

\subsection{KV Cache and Attention Memory Efficiency}

The KV cache stores key and value projections for all previous tokens, growing linearly with sequence length and consuming a large fraction of GPU memory during decoding~\cite{zhang2023h2o}.
H$_2$O~\cite{zhang2023h2o} proposes a dynamic eviction policy that retains only a small subset of ``heavy hitter'' tokens, reducing KV cache memory footprint while preserving model quality.
FlashAttention~\cite{dao2022flashattention} takes a complementary approach: rather than managing cache size, it restructures the attention computation to minimize traffic between GPU high-bandwidth memory and on-chip SRAM via tiling, yielding significant memory savings and runtime speedup without approximation.
Both lines of work assume that the underlying hardware executes attention kernels predictably as a function of tensor size.
Our results challenge this assumption for the MPS backend: we observe that KV cache growth causes latency to become non-monotonic in a way that FlashAttention-style IO analysis cannot predict, since the MPS backend does not expose the same memory hierarchy controls as CUDA.

\subsection{On-Device and Apple Silicon Inference}

Inference on devices with constrained memory has received growing attention.
Alizadeh et al.~\cite{alizadeh2024llmflash} study LLM inference when model parameters exceed available DRAM on Apple Silicon, introducing techniques to stream weights from flash memory to unified memory efficiently.
That work focuses on parameter loading bandwidth rather than execution dynamics during decoding and does not investigate KV cache behavior.
Rajesh et al.~\cite{rajesh2025production} compare five Apple Silicon runtimes (MLX, MLC-LLM, Ollama, llama.cpp, and PyTorch MPS) across Qwen-2.5 models at prompt lengths up to 100,000 tokens, reporting throughput, TTFT, and long-context behavior.
They find that PyTorch MPS is limited on large models and long contexts, but characterize performance as monotonically degrading rather than exhibiting discrete, non-monotonic instability.
Benazir and Lin~\cite{benazir2025profiling} profile Apple Silicon's unified memory architecture from a quantization perspective, finding it cost-effective for large models; however, their analysis focuses on static benchmarks and quantization overhead rather than decode-phase dynamics.
Horton et al.~\cite{horton2024kvprediction} introduce KV Prediction to reduce time-to-first-token on edge devices by amortizing the KV cache construction cost; their work assumes KV cache is universally beneficial, and our results provide complementary evidence that KV cache growth can itself become a source of instability during decoding on the MPS backend.
The PyTorch MPS backend~\cite{pytorch2022mps} enables GPU-accelerated inference on Apple Silicon via the Metal Performance Shaders framework~\cite{apple2026mps}, but its kernel selection and scheduling logic is not publicly documented, making it difficult to predict or diagnose performance anomalies from first principles.
Our work is the first to empirically characterize pronounced non-monotonic latency behavior in MPS autoregressive decoding, quantify a KV-cache speedup collapse at pathological lengths, and characterize the asymmetric onset/recovery pattern through controlled ablations.

\section{Experimental Setup}

\subsection{Experimental Configuration}

Experiments are conducted on a MacBook Pro equipped with an Apple M3 Max processor (14-core CPU, 30-core GPU, 36~GB unified memory) running macOS 14.8.1 (Sonoma). The software stack includes Python~3.9.6, PyTorch~2.8.0~\cite{paszke2019pytorch}, HuggingFace Transformers~4.57.6~\cite{wolf2020transformers}, and LitServe~0.2.16~\cite{lightningai2026litserve}, using the MPS backend for GPU-accelerated execution.

We evaluate four models from the GPT-2 family: DistilGPT-2~\cite{sanh2019distilbert}, GPT-2~\cite{radford2019gpt2}, GPT-2 Medium~\cite{radford2019gpt2}, and GPT-2 Large~\cite{radford2019gpt2}. Inference is tested on two execution devices: CPU and MPS (GPU via Metal Performance Shaders).

\subsection{Benchmark Design}

All experiments are served through LitServe~\cite{lightningai2026litserve}, a production-grade Python inference serving framework.
Using a real serving stack rather than bare model calls ensures that overheads from request handling, worker scheduling, and response serialization are present throughout, making the observed anomalies relevant to actual deployment scenarios.
Crucially, the anomalies manifest at the decode phase (token-level timing shows prefill is negligible), which is decoupled from LitServe's request handling, indicating the instability originates in the MPS backend's kernel execution rather than framework scheduling.

Each experiment runs in an isolated server process with:
\begin{itemize}
\item 3 warmup runs (discarded),
\item 5 measured runs (averaged),
\item deterministic decoding (fixed seed),
\item token-level timing for selected experiments.
\end{itemize}

Generation lengths range from 128 to 768 tokens.

\paragraph{Measurement methodology.}
Baseline scaling experiments (Section~\ref{sec:baseline}) measure end-to-end request latency without per-step instrumentation, representing natural execution conditions.
For per-token profiling experiments (KV cache ablation and instability probe), a device synchronization barrier is inserted after each decode step to obtain accurate step-level timing.
This barrier forces the GPU to fully flush outstanding work before the timer advances, which is necessary for correctness but introduces serialization overhead absent in the uninstrumented path.
Consequently, absolute latency values from per-token profiling experiments are not directly comparable to baseline scaling results.
Additionally, because 512 tokens falls near the onset threshold, latency at that specific length is highly variable across independent runs (standard deviation 8--14~s); the values reported in each section are the mean of the 5 measured runs for that experiment and will therefore differ slightly across sections.
All comparisons between KV cache enabled and disabled conditions use the same instrumentation, making them internally consistent for relative analysis.
Importantly, all reported non-monotonic behavior is present in uninstrumented end-to-end measurements (Section~\ref{sec:baseline}); per-step instrumentation is used only for attribution of latency to prefill vs.~decode, not as the basis for the anomaly observation itself.

\section{Results}

\subsection{Baseline Scaling and CPU--MPS Crossover}
\label{sec:baseline}

At shorter generation lengths (128--256 tokens), latency scales approximately linearly with output length across all models and modes. We observe a clear model-size-dependent CPU--MPS crossover, summarized in Table~\ref{tab:baseline_128}.

\begin{table}[h]
\centering
\caption{Average latency (s) at 128 generated tokens across models and devices. MPS speedup is relative to CPU.}
\label{tab:baseline_128}
\begin{tabular}{lrrc}
\toprule
Model & CPU & MPS & MPS speedup \\
\midrule
DistilGPT-2  & 0.84 & 0.78 & $1.1\times$ \\
GPT-2        & 1.38 & 0.99 & $1.4\times$ \\
GPT-2 Medium & 4.01 & 1.38 & $2.9\times$ \\
GPT-2 Large  & 7.13 & 2.21 & $3.2\times$ \\
\bottomrule
\end{tabular}
\end{table}

This crossover validates the experimental setup and confirms that MPS advantage grows with model size at short generation lengths. This scaling behavior is visualized across all models in Figure~\ref{fig:latency_scaling}.

\begin{figure}[t]
  \centering
  \includegraphics[width=\linewidth]{fig1_latency_scaling.pdf}
  \caption{Latency scaling by model and device. Each panel shows average latency vs.\ generated tokens for one model. Non-monotonic spikes are visible for GPT-2 Medium and GPT-2 Large on MPS at long sequence lengths, with both models entering an anomalous execution regime beyond model-size-dependent token thresholds.}
  \label{fig:latency_scaling}
\end{figure}

\subsection{Emergence of Non-Monotonic Scaling}

At longer generation lengths, the expected smooth scaling breaks down on the MPS backend.
The anomalies are model-size-dependent and occur at different token thresholds, as summarized in Table~\ref{tab:anomaly_summary}.

\begin{table}[h]
\centering
\caption{Non-monotonic latency anomalies on MPS. Spike latency is the average at the anomalous token count. Reference latency is the nearest non-anomalous configuration (immediately below the spike). Slowdown is spike / reference.}
\label{tab:anomaly_summary}
\begin{tabular}{llrrrr}
\toprule
Model & Mode & Anomalous & Spike & Reference & Slowdown \\
      &      & token     & (s)   & (s)       &          \\
\midrule
GPT-2 Medium & MPS & 512 & 84.5  & 5.1 at 384  & $16.6\times$ \\
GPT-2 Large  & MPS & 384 & 103.4 & 4.9 at 256  & $21.1\times$ \\
\bottomrule
\end{tabular}
\end{table}

These effects are strictly non-monotonic: the anomalous configuration performs worse than both shorter and longer generation lengths.
Smaller models (DistilGPT-2, GPT-2) do not exhibit the anomaly within the 128--768 token range tested, suggesting a model-size threshold for the instability to manifest.
Latency slowdown reaches up to $21\times$ in the worst case (GPT-2 Large / MPS at 384 tokens).

\subsection{Instability Probe}

A focused instability sweep over 480--656 tokens at 16-token spacing reveals a broad anomalous regime with sharp boundaries.
Latency is in the normal range at 480--496 tokens (9.2--9.5s), then jumps abruptly at 512 tokens (89.2s).
The elevated regime persists through 624 tokens (83.2--110.1s across the sampled points), and then collapses back to normal at 640--656 tokens (8.8s and 8.2s, respectively).

These measurements show that the instability is not an isolated single-token outlier at 512, but a distinct high-latency execution regime.
Across the 5 measured runs, the anomalous regime exhibits consistent high latency (standard deviation 8--14~s at 512 tokens), indicating stable but degraded execution rather than transient noise. The anomalous regime is consistently observed across independent runs and is not attributable to measurement variation.
The regime transition occurs within one 16-token step on both sides (onset between 496 and 512, recovery between 624 and 640), which is inconsistent with smooth scaling and supports a discrete backend execution-path switch. These sharp boundaries are evident in Figure~\ref{fig:instability_probe}.
Per-step timing in the instrumented runs confirms that prefill cost is negligible throughout: at 512 tokens the prefill phase completes in 0.11~s (0.1\% of total latency), and at the adjacent normal point (496 tokens) prefill is 0.08~s (0.8\% of total latency).
The latency difference between 496 tokens (9.5~s) and 512 tokens (89.2~s) is therefore attributable entirely to the decode phase (9.4~s vs 89.1~s respectively), not to any change in prefill cost.
Table~\ref{tab:instability_probe} reports the full sweep.

\begin{table}[h]
\centering
\caption{Full instability probe for GPT-2 Medium on MPS, 16-token intervals from 480 to 656. The anomalous regime (512--624 tokens) is set in bold; boundary transitions occur within a single 16-token step on each side.}
\label{tab:instability_probe}
\begin{tabular}{rrl}
\toprule
\texttt{max\_tokens} & Latency (s) & Regime \\
\midrule
480 & 9.2   & normal \\
496 & 9.5   & normal \\
\textbf{512} & \textbf{89.2}  & \textbf{anomalous} \\
\textbf{528} & \textbf{83.2}  & \textbf{anomalous} \\
\textbf{544} & \textbf{97.0}  & \textbf{anomalous} \\
\textbf{560} & \textbf{110.1} & \textbf{anomalous} \\
\textbf{576} & \textbf{99.9}  & \textbf{anomalous} \\
\textbf{592} & \textbf{99.2}  & \textbf{anomalous} \\
\textbf{608} & \textbf{98.9}  & \textbf{anomalous} \\
\textbf{624} & \textbf{94.4}  & \textbf{anomalous} \\
640 & 8.8   & normal \\
656 & 8.2   & normal \\
\bottomrule
\end{tabular}
\end{table}

\begin{figure}[t]
  \centering
  \includegraphics[width=0.85\linewidth]{fig2_instability_probe.pdf}
  \caption{Instability probe: average latency vs.\ \texttt{max\_tokens} for GPT-2 Medium on MPS at 16-token intervals. The shaded region marks the anomalous regime (512--624 tokens). Dashed lines indicate the onset and recovery boundaries.}
  \label{fig:instability_probe}
\end{figure}

\subsection{KV Cache Ablation}
\label{sec:kvcache}

We perform a controlled ablation on GPT-2 Medium using MPS, varying KV cache state across 128--768 tokens.
All other conditions (model, prompt, seed, instrumentation) are held constant.
The KV-cache-disabled condition uses the same protocol as all other experiments (3 warmup runs and 5 measured runs); the $O(n^2)$ overhead at long sequence lengths is elevated but tractable within this run count, and is sufficient to test for the presence or absence of the anomaly.

With KV cache enabled, the anomaly at 512 tokens is extreme: latency spikes to 83.2s, compared to 5.1s at 384 tokens and 8.8s at 640 tokens.
With KV cache disabled in the rerun, non-monotonicity remains: cache-off latency is 161.6s at 512 tokens and 139.0s at 640 tokens.
At the anomalous 512-token point, cache-off is only $1.9\times$ slower (161.6s vs 83.2s), indicating that cache-enabled execution is unusually inefficient there relative to its surrounding points.
At this point, cache-off execution, despite its higher theoretical complexity, approaches cache-on latency, directly contradicting the expected computational advantage of KV caching. In this regime, backend execution effects dominate algorithmic complexity to the extent that an $O(n)$ method exhibits near-$O(n^2)$ practical performance, and the practical benefit of KV caching is largely neutralized.

Variance at longer cache-off lengths is substantial and should be interpreted explicitly: standard deviation is 1.1s at 512, 21.5s at 640, and 40.2s at 768, with observed run-level ranges of 159.9--162.7s (512), 121.2--175.1s (640), and 168.7--276.1s (768). Accordingly, upper-tail metrics (p95 and max) increase monotonically with token length in this rerun, while central tendency (mean and median) still shows a 512 $>$ 640 inversion.

These results show a clear KV-cache trade-off on MPS: cache-enabled decoding remains nominally faster at all tested lengths, but its expected speedup fails to materialize at pathological lengths, collapsing from order-of-magnitude gains to near parity rather than degrading smoothly.
Because all experiments use deterministic greedy decoding with a fixed random seed, generated outputs are identical across KV cache enabled and disabled conditions, confirming that the ablation affects only execution dynamics, not output quality.

Notably, disabling KV cache reintroduces the theoretically less efficient $O(n^2)$ attention computation, and while it reduces the cache-on spike contrast, it does not yield strictly monotonic execution in this run. This suggests that the anomalous behavior is consistent with backend-dependent execution path changes in the MPS backend, where cache growth interacts with backend dispatch decisions to produce localized speedup collapse. Importantly, the anomalous regime persists under both KV-cache-enabled and KV-cache-disabled execution, indicating that KV caching does not cause the anomaly but modulates its severity. The comparison is shown in Figure~\ref{fig:kv_cache_ablation}.

Table~\ref{tab:kv_ablation} reports the latency comparison across 128--768 tokens. At non-anomalous points, cache-off / cache-on is large ($4.9\times$--$20.3\times$). At the anomalous 512-token point the ratio collapses to $1.9\times$: cache-off is only modestly slower despite doing substantially more arithmetic, consistent with the spike under cache-on being a hardware-regime artifact rather than genuine algorithmic work.

\begin{table}[h]
\centering
\caption{KV cache ablation for GPT-2 Medium on MPS: latency (s) with cache enabled vs.\ disabled. The \textbf{bold row} marks the anomalous configuration (512 tokens). The ratio (cache-off / cache-on) collapses at that point, showing the spike is not due to excess computation.}
\label{tab:kv_ablation}
\begin{tabular}{r rr r}
\toprule
Tokens & Cache ON (s) & Cache OFF (s) & Ratio \\
\midrule
128 & 1.4  & 6.8   & $4.9\times$ \\
256 & 2.9  & 20.6  & $7.1\times$ \\
384 & 5.1  & 42.7  & $8.4\times$ \\
\textbf{512} & \textbf{83.2} & 161.6 & $1.9\times$ \\
640 & 8.8  & 139.0 & $15.8\times$ \\
768 & 9.7  & 196.6 & $20.3\times$ \\
\bottomrule
\end{tabular}
\end{table}

\paragraph{Memory pressure as a confound.}
To rule out memory exhaustion as a driver of instability, we profile peak MPS-allocated memory across all token lengths using a dedicated memory-instrumented run (kept separate from latency measurements to avoid sampling overhead).
For GPT-2 Medium on MPS, peak allocated memory increases monotonically and uniformly across all token lengths, as shown in Table~\ref{tab:memory_profile}: from 1426\ MB at 128 tokens to 1550\ MB at 768 tokens, in increments of approximately 25\ MB per 128-token step, with no discontinuity at the anomalous 512-token point.

\begin{table}[h]
\centering
\caption{Peak MPS-allocated memory for GPT-2 Medium across token lengths. The per-step increment is approximately 25\ MB throughout, with no discontinuity at the 512-token anomaly boundary.}
\label{tab:memory_profile}
\begin{tabular}{r r r}
\toprule
Tokens & Peak memory (MB) & Increment (MB) \\
\midrule
128 & 1426 & --- \\
256 & 1451 & $+25$ \\
384 & 1476 & $+25$ \\
512 & 1500 & $+24$ \\
640 & 1525 & $+25$ \\
768 & 1550 & $+25$ \\
\bottomrule
\end{tabular}
\end{table}

The same smooth monotonic pattern holds for GPT-2 Large.
Because peak memory at the latency spike (512 tokens, 93.2s in the memory-profile run) is indistinguishable from the surrounding normal configurations, the anomalous latency behavior is unlikely to be attributed to memory pressure or allocation failures, and is more consistent with differences in kernel execution paths within the MPS backend.

\begin{figure}[t]
  \centering
  \includegraphics[width=\linewidth]{fig3_kv_cache_ablation.pdf}
  \caption{KV cache ablation for GPT-2 Medium on MPS. With KV cache enabled, latency exhibits a sharp non-monotonic spike at 512 tokens. With KV cache disabled, latency is higher throughout and remains non-monotonic. Together, these results indicate a KV-cache trade-off: strong speedups at most lengths, but severe speedup collapse at specific pathological lengths on MPS.}
  \label{fig:kv_cache_ablation}
\end{figure}

\subsection{Prompt-Length Analysis}
\label{sec:prompt}

To test whether the instability is governed by the generation budget alone (\texttt{max\_tokens}) or by the total sequence length (prompt tokens + generated tokens), we hold \texttt{max\_tokens} constant and vary the input prompt length across two conditions: a short prompt of 3 tokens (\textit{``Hello world.''}) and a long prompt of approximately 65 tokens (a technical passage).
At each \texttt{max\_tokens} setting the corresponding total sequence lengths are therefore offset by roughly 62 tokens.

The results reveal an onset asymmetry: at \texttt{max\_tokens}~=~496, both prompts are in the normal regime (7.1s and 9.5s respectively), yet at \texttt{max\_tokens}~=~512 the short prompt remains normal (7.3s) while the long prompt enters the anomalous regime (83.9s).
Under a generation-budget-only hypothesis both prompts would be affected equally, since they share the same \texttt{max\_tokens} value.
The observed divergence at 512 tokens (where the long prompt's total sequence length ($\approx$577) crosses the instability threshold while the short prompt's total ($\approx$515) does not) is consistent with total sequence length playing a role in onset.

The recovery boundary, however, does not show a corresponding asymmetry: at \texttt{max\_tokens}~=~624 both prompts are slow (104.2s and 92.0s), and at \texttt{max\_tokens}~=~640 both recover (9.5s and 8.6s).
Notably, at \texttt{max\_tokens}~=~624 the short prompt (total $\approx$627 tokens) is actually slower than the long prompt (total $\approx$689 tokens) despite having a lower total sequence length, indicating that total sequence length alone does not monotonically predict latency within the anomalous regime.
Because the short prompt's total sequence length at \texttt{max\_tokens}~=~640 ($\approx$643) is well below the long prompt's ($\approx$705), yet both recover to similar latencies, recovery cannot be driven by total sequence length.
Instead, recovery correlates strongly with \texttt{max\_tokens}: both prompts recover at 640 regardless of their different total lengths.

The onset pattern is less clear: the long prompt enters the anomalous regime at max\_tokens~=~512 (total $\approx$577), while the short prompt remains normal until at least max\_tokens~=~624 (total $\approx$627), suggesting total sequence length influences onset. However, the non-monotonic behavior within the anomalous regime at max\_tokens~=~624 indicates the underlying mechanism is more complex than a simple threshold. These results suggest that both \texttt{max\_tokens} and total sequence length play roles, but their interaction is not straightforward. This divergence is illustrated in Figure~\ref{fig:prompt_ablation}.
Table~\ref{tab:prompt_ablation} gives the exact latencies and total sequence lengths for each condition.

\begin{table}[h]
\centering
\caption{Prompt-length ablation results for GPT-2 Medium on MPS. Short prompt: 3 tokens; long prompt: 65 tokens. Total sequence length = \texttt{max\_tokens} + prompt tokens. At \texttt{max\_tokens}~=~512 the prompts diverge: the long prompt's total sequence length ($\approx$577) crosses the instability threshold while the short prompt's ($\approx$515) does not. Both recover simultaneously at 640, despite different total lengths.}
\label{tab:prompt_ablation}
\begin{tabular}{r rr rr}
\toprule
& \multicolumn{2}{c}{Short prompt (3 tok)} & \multicolumn{2}{c}{Long prompt (65 tok)} \\
\cmidrule(lr){2-3}\cmidrule(lr){4-5}
\texttt{max\_tokens} & Total seq & Latency (s) & Total seq & Latency (s) \\
\midrule
496 & $\approx$499 & 7.1  & $\approx$561 & 9.5  \\
512 & $\approx$515 & 7.3  & $\approx$577 & 83.9 \\
624 & $\approx$627 & 104.2 & $\approx$689 & 92.0 \\
640 & $\approx$643 & 9.5  & $\approx$705 & 8.6  \\
\bottomrule
\end{tabular}
\end{table}

\begin{figure}[t]
  \centering
  \includegraphics[width=0.75\linewidth]{fig4_prompt_ablation.pdf}
  \caption{Prompt-length ablation for GPT-2 Medium on MPS. At \texttt{max\_tokens}~=~512 the long prompt (total $\approx$577 tokens) enters the anomalous execution regime while the short prompt (total $\approx$515 tokens) remains normal, consistent with total sequence length influencing onset.}
  \label{fig:prompt_ablation}
\end{figure}

\section{Discussion}

\subsection{Discrete Execution Regimes}

The instability probe (Table~\ref{tab:instability_probe}) reveals that the transition between normal and anomalous execution is discrete, not gradual: between 496 and 512 tokens (a single 16-token step), latency jumps from 9.5s to 89.2s with no intermediate states observed.
This is inconsistent with any smooth algorithmic cost model and is consistent with discrete kernel selection within the MPS backend, though the exact mechanism cannot be verified without instrumentation.
Similar threshold behavior is well documented in GPU computing when tensor dimensions cross tiling boundaries or when memory access patterns shift between coalesced and non-coalesced regimes, but the MPS backend does not expose the instrumentation needed to confirm the exact trigger.
We emphasize that our results characterize observable execution behavior; we do not have visibility into Metal kernel dispatch or internal scheduling decisions, and therefore interpret the anomaly as evidence of a regime shift rather than a confirmed implementation-level change.
The model-size dependence in Table~\ref{tab:anomaly_summary} is consistent with this hypothesis: larger models have wider hidden dimensions, so the critical total tensor size (hidden dimension $\times$ sequence length) is reached at shorter token counts.
DistilGPT-2 and GPT-2 do not exhibit the anomaly within the 128--768 token range, while GPT-2 Medium and GPT-2 Large do, with the larger model crossing the threshold earlier.

\subsection{KV Cache Under Regime Shifts}

During autoregressive decoding with KV cache enabled, the key-value tensors grow by one row per generated token, creating a monotonically increasing state that is read and written at every decode step.
The KV cache ablation (Table~\ref{tab:kv_ablation}) shows that removing this growing state does not eliminate instability entirely: cache-off latency remains non-monotonic in this run (e.g., 512 tokens slower than 640), even as the sharp cache-on spike contrast is reduced.
Because cache-off removes the growing KV tensor yet residual non-monotonicity remains, the anomaly cannot be attributed to KV growth alone; both modes exhibit non-monotonic behavior, but the magnitude and location of instability differ markedly (83.2s spike at 512 tokens with cache-on versus a smaller 512-over-640 inversion with cache-off).
The diagnostic signature is in the cache-off / cache-on latency ratio: at non-anomalous points this ratio is large ($4.9\times$--$20.3\times$), confirming that KV cache is strongly beneficial there.
At the anomalous 512-token point the ratio collapses to $1.9\times$, meaning that with cache disabled the system is only modestly slower despite doing substantially more arithmetic.
This inversion represents a breakdown of the expected correspondence between algorithmic complexity and runtime performance; it does not imply algorithmic reversal, but rather that constant-factor effects and backend execution overhead dominate theoretical complexity in this regime. We therefore characterize KV caching on MPS as \emph{regime-sensitive}: it delivers strong benefits within stable execution regions, but its practical advantage can collapse sharply when the backend shifts execution regime.

\subsection{Onset and Recovery Patterns}

The onset and recovery boundaries show different correlations, as the prompt-length ablation (Section~\ref{sec:prompt} and Table~\ref{tab:prompt_ablation}) reveals.
Onset correlates with total sequence length but does not follow a single fixed threshold: at \texttt{max\_tokens}~=~512, the long prompt (total $\approx$577 tokens) enters the anomalous regime while the short prompt (total $\approx$515 tokens) does not; however, the short prompt enters at \texttt{max\_tokens}~=~624 (total $\approx$627 tokens), suggesting total sequence length is a necessary but not sufficient condition for onset. These observations indicate a more complex interaction between sequence length and backend execution state.
Recovery, however, correlates with \texttt{max\_tokens} alone: both prompts recover simultaneously at \texttt{max\_tokens}~=~640 regardless of their different total sequence lengths.
This asymmetry is consistent with two different empirical dependencies: onset associated with actual KV cache growth during decoding (total sequence length), and recovery associated with the static \texttt{max\_tokens} budget set at initialization.
Such a mixed static/dynamic dispatch is plausible in backends that pre-plan buffer layouts at graph compilation time while handling per-step tensor growth at runtime. The interaction between these two mechanisms exhibits an asymmetry: entry and exit conditions differ in a manner suggestive of hysteresis-like behavior, creating a bounded but unstable regime. We use 'hysteresis-like' descriptively to indicate asymmetric entry and exit conditions, without implying a specific physical mechanism.

\subsection{Implications for Practice}

These findings have direct consequences for practitioners deploying autoregressive models on Apple Silicon.
First, KV cache cannot be assumed to deliver stable speedup: at certain sequence lengths on MPS, latency exhibits sharp spikes up to $21\times$ above the pre-threshold baseline (Table~\ref{tab:anomaly_summary}), collapsing the typical caching advantage to as little as $2\times$ at those lengths.
Second, benchmark methodology matters: coarse throughput measurements at short sequence lengths will miss these effects entirely, as the instability only manifests beyond model-size-dependent token thresholds.
Reporting performance at a single generation length, or only at 128--256 tokens, gives a misleading picture of real-world behavior for long-context workloads.
Third, because the recovery boundary is governed by \texttt{max\_tokens} rather than actual output length, it is in principle adjustable at deployment time: reducing the generation budget below the recovery threshold restores normal performance, at the cost of a shorter maximum output.
For applications where context length can be bounded, this is a practical mitigation that requires no changes to the model or backend.

\subsection{Limitations}

All experiments were conducted on a single machine (MacBook Pro, Apple M3 Max) running a specific software configuration (PyTorch~2.8.0, macOS~14.8.1). We cannot rule out that the observed behavior is specific to this chip variant, driver version, or PyTorch release. Reproducing on additional Apple Silicon generations (M1, M2) or PyTorch versions would strengthen the generalizability of the findings; we leave this to future work. Additionally, because the MPS backend does not expose kernel-level instrumentation, the proposed execution-regime hypothesis (while consistent with all observations) cannot be directly verified without access to Metal shader internals.

\subsection{Takeaway}

\begin{quote}
\textbf{On the MPS backend, autoregressive inference exhibits abrupt, threshold-like latency changes consistent with an underlying backend-dependent shift in execution behavior, where these effects can dominate algorithmic complexity, causing performance to degrade sharply at specific sequence lengths and violating standard assumptions of smooth scaling.}
\end{quote}

\section{Conclusion}

This study characterizes pronounced non-monotonic latency behavior in autoregressive inference on Apple MPS, and shows that KV cache introduces a practical trade-off: large speedups at most lengths but sharp speedup collapse at specific pathological lengths. We also characterize an asymmetric onset/recovery pattern under cache-enabled decoding: onset correlates with total sequence length (prompt + generated tokens), while recovery correlates with the static \texttt{max\_tokens} budget set at initialization. The instability manifests only at model-size-dependent token thresholds: GPT-2 Medium and GPT-2 Large are affected within the 128--768 token range while smaller models are not, consistent with dependence on model capacity and KV state size. Disabling KV cache changes the latency profile and reduces the sharp cache-on spike contrast, but does not fully remove non-monotonic behavior in this measurement set. Practically, tuning \texttt{max\_tokens} away from observed anomalous regions is a zero-model-change mitigation, and coarse benchmarks at short sequence lengths can systematically miss this behavior. These findings underscore the importance of hardware-aware, sequence-length-resolved evaluation for any deployment targeting long-context generation on Apple Silicon.

\section*{Code Availability}
The benchmark harness and configuration files used in this study are available at \url{https://github.com/willyfh/kv-cache-mps-instability}.

\bibliographystyle{unsrt}
\bibliography{references}

\end{document}